\section{Experiment 2: Controlled Defect Injection}
\label{sec:results-exp2}

\subsection{Detection on the Structural Bands}
\label{sec:results-exp2-structural}

Figure~\ref{fig:exp2-detection} shows detection rates on the 28 structural
injections, where the consequence lies in a different file from the edit.%
\footnote{Numbers reflect the corrected two-stage adjudication described in
Section~\ref{sec:injection-design}, scored by the five-judge panel. Exact
rates and Wilson intervals are in
Table~\ref{tab:exp2-detection} (Appendix).}

\begin{figure}[htbp]
    \centering
    % !TeX root = ../thesis.tex
% GENERATED FILE - do not edit by hand; your changes will be lost.
% Produced by scripts/generate_thesis_figures.py
% Regenerate: python3 scripts/generate_thesis_figures.py
% Computed from:
%   experiments/2026-09-23_five_judge_panel/RESULTS.json (Experiment 2 detection)
%
\begin{tikzpicture}
\node[anchor=west,font=\footnotesize] at (-4.4000,0.0000) {\baseline{} (diff only)};
\draw[black!70,line width=0.5pt] (0.0000,0.0000) -- (0.7239,0.0000);
\draw[black!70,line width=0.5pt] (0.0000,-0.0550) -- (0.0000,0.0550);
\draw[black!70,line width=0.5pt] (0.7239,-0.0550) -- (0.7239,0.0550);
\draw[fill=white,draw=black!75,line width=0.6pt] (6.0000,0.0000) circle (2.100pt);
\node[anchor=east,font=\scriptsize] at (7.5000,0.0000) {0/28};
\node[anchor=west,font=\footnotesize] at (-4.4000,-0.5000) {\rag{}};
\fill[ragorange] (0.0000,-0.6450) rectangle (1.0714,-0.3550);
\draw[black!70,line width=0.5pt] (0.4727,-0.5000) -- (2.1355,-0.5000);
\draw[black!70,line width=0.5pt] (0.4727,-0.5550) -- (0.4727,-0.4450);
\draw[black!70,line width=0.5pt] (2.1355,-0.5550) -- (2.1355,-0.4450);
\draw[fill=white,draw=black!75,line width=0.6pt] (6.0000,-0.5000) circle (2.100pt);
\node[anchor=east,font=\scriptsize] at (7.5000,-0.5000) {5/28};
\node[anchor=west,font=\footnotesize] at (-4.4000,-1.0000) {\hybrid{}};
\fill[hybridgreen] (0.0000,-1.1450) rectangle (4.5000,-0.8550);
\draw[black!70,line width=0.5pt] (3.3987,-1.0000) -- (5.2394,-1.0000);
\draw[black!70,line width=0.5pt] (3.3987,-1.0550) -- (3.3987,-0.9450);
\draw[black!70,line width=0.5pt] (5.2394,-1.0550) -- (5.2394,-0.9450);
\draw[fill=white,draw=black!75,line width=0.6pt] (5.5000,-1.0000) circle (2.100pt);
\node[anchor=east,font=\scriptsize] at (7.5000,-1.0000) {21/28};
\node[anchor=west,font=\footnotesize] at (-4.4000,-1.5000) {\kg{} (deployed)};
\fill[kgblue] (0.0000,-1.6450) rectangle (3.8571,-1.3550);
\draw[black!70,line width=0.5pt] (2.7498,-1.5000) -- (4.7577,-1.5000);
\draw[black!70,line width=0.5pt] (2.7498,-1.5550) -- (2.7498,-1.4450);
\draw[black!70,line width=0.5pt] (4.7577,-1.5550) -- (4.7577,-1.4450);
\draw[fill=white,draw=black!75,line width=0.6pt] (6.0000,-1.5000) circle (2.100pt);
\node[anchor=east,font=\scriptsize] at (7.5000,-1.5000) {18/28};
\node[anchor=west,font=\footnotesize] at (-4.4000,-2.0000) {\kg{} + inheritance edges};
\fill[kgblue] (0.0000,-2.1450) rectangle (5.5714,-1.8550);
\draw[black!70,line width=0.5pt] (4.6413,-2.0000) -- (5.8811,-2.0000);
\draw[black!70,line width=0.5pt] (4.6413,-2.0550) -- (4.6413,-1.9450);
\draw[black!70,line width=0.5pt] (5.8811,-2.0550) -- (5.8811,-1.9450);
\draw[fill=white,draw=black!75,line width=0.6pt] (6.0000,-2.0000) circle (2.100pt);
\node[anchor=east,font=\scriptsize] at (7.5000,-2.0000) {26/28};
\node[anchor=west,font=\footnotesize] at (-4.4000,-2.5000) {\kg{} idealised (ceiling)};
\fill[kgblue] (0.0000,-2.6450) rectangle (5.5714,-2.3550);
\draw[black!70,line width=0.5pt] (4.6413,-2.5000) -- (5.8811,-2.5000);
\draw[black!70,line width=0.5pt] (4.6413,-2.5550) -- (4.6413,-2.4450);
\draw[black!70,line width=0.5pt] (5.8811,-2.5550) -- (5.8811,-2.4450);
\draw[fill=white,draw=black!75,line width=0.6pt] (6.0000,-2.5000) circle (2.100pt);
\node[anchor=east,font=\scriptsize] at (7.5000,-2.5000) {26/28};
\node[anchor=east,font=\scriptsize] at (7.5000,0.4000) {detected};
\draw[thin,black!55] (0.0000,-2.7750) -- (6.0000,-2.7750);
\draw[thin,black!55] (0.0000,-2.7750) -- (0.0000,-2.8550);
\node[anchor=center,font=\scriptsize] at (0.0000,-3.0550) {0};
\draw[thin,black!55] (1.2000,-2.7750) -- (1.2000,-2.8550);
\node[anchor=center,font=\scriptsize] at (1.2000,-3.0550) {0.2};
\draw[thin,black!55] (2.4000,-2.7750) -- (2.4000,-2.8550);
\node[anchor=center,font=\scriptsize] at (2.4000,-3.0550) {0.4};
\draw[thin,black!55] (3.6000,-2.7750) -- (3.6000,-2.8550);
\node[anchor=center,font=\scriptsize] at (3.6000,-3.0550) {0.6};
\draw[thin,black!55] (4.8000,-2.7750) -- (4.8000,-2.8550);
\node[anchor=center,font=\scriptsize] at (4.8000,-3.0550) {0.8};
\draw[thin,black!55] (6.0000,-2.7750) -- (6.0000,-2.8550);
\node[anchor=center,font=\scriptsize] at (6.0000,-3.0550) {1};
\node[anchor=center,font=\footnotesize] at (3.0000,-3.4950) {Detection rate};
\fill[black!70] (0.0000,-4.3910) rectangle (0.4200,-4.1590);
\node[anchor=west,font=\scriptsize] at (0.5500,-4.2750) {structural injections, consequence outside the diff ($n = 28$), with 95\,\% CI};
\draw[fill=white,draw=black!75,line width=0.6pt] (0.2100,-4.6750) circle (2.100pt);
\node[anchor=west,font=\scriptsize] at (0.5500,-4.6750) {local control, consequence inside the diff ($n = 12$)};
\end{tikzpicture}

    \caption{Detection rate by arm on both bands. Bars are structural
             injections with Wilson 95\,\% intervals; open markers are local
             controls. Structural-band rates increase from zero to
             near-ceiling, while local-control rates stay at or near
             ceiling.}
    \label{fig:exp2-detection}
\end{figure}

The baseline detects nothing (zero of 28) while the deployed graph arm
detects eighteen. The exact McNemar contrast is $+0.64$, $p < 0.0001$,
clearing the pre-registered threshold of $+0.30$ at $p < 0.05$.

Retrieval detects five of 28. On pull requests it had the best total coverage;
on defects whose
consequences sit outside the diff it remains substantially below the graph.
Embedding similarity retrieves code that \emph{resembles} the change, but a
file that breaks because it calls a renamed function need not resemble it.
The two mechanisms are not competing implementations of the same idea.

\hybrid{} detects twenty-one, numerically above the graph arm. The paired
difference does not provide evidence that the hybrid outperforms the graph
arm.

\subsection{The Local Control Bands}
\label{sec:results-exp2-control}

On the twelve local injections, whose consequences are visible in the diff,
all arms detect at or near ceiling (the hybrid arm misses one local case;
every other arm scores $12/12$). The graph block's advantage appears exactly
where structural reasoning is required and vanishes where it is not.

The graph block contains the dependents' filenames, so naming one is not
by itself proof of reasoning. Judges must confirm that the review links
the named file to the change, and filtering the block to true-dependent
edges raises detection further (Appendix~\ref{app:supporting-components}),
so what is named matters. A block of the same size with edges to unrelated
symbols yields 1 of 28, so the detections do not come from memory of these
repositories.

\subsection{Effect of Inheritance Edges}
\label{sec:results-exp2-inheritance}

The deployed builder carries no inheritance edges. Adding them raises
detection from 18 to 26 of 28, matching the idealised resolver. With the model
held fixed, this attributes the eight
recovered cases to relations omitted by the deployed builder. The two
remaining misses persist under the idealised resolver and cannot be assigned
to graph construction alone: one review identifies affected classes without
satisfying the exact filename gate, while the other remains generic despite
complete evidence. Perfect structural evidence therefore does not guarantee a
specific review comment. Appendix~\ref{app:supporting-exp2-misses} analyses
both cases; Appendix~\ref{app:supporting-exp2-contrasts} reports the secondary
paired contrasts.

\subsection{Supporting Mechanism Probes}

Component probes isolate dependency lists, resolved call edges, test-file
links, and matched-volume noise. Across their respective configurations they
support the same structural interpretation: relevant relationships drive
detection, while unrelated context does not. Appendix
\ref{app:supporting-components} reports the protocols and statistics.
